\section{The ordered spectral law and its transfer}\label{sec:multiple}

\subsection{The surviving model harmonic space}
Fix $R,A$. For $k\in K$, let $\eta_k^R$ be its central harmonic
representative. Its harmonic link coefficients vanish. The exact-cone
expansion \eqref{eq:cone-harmonic-expansion} gives an end primitive
$\beta_k$ of size $O(r^{\nu-1})$, for some $\nu>0$, with differentiated
bounds. If $\chi_\rho$ equals one near each tip and zero above $2\rho$,
then replacing the form on its ends by
\begin{equation}\label{eq:closed-compact-approximation}
 \gamma_{k,\rho}=\eta_k^R-d(\chi_\rho\beta_k)
\end{equation}
gives a compactly supported closed representative of $k$ converging in
$L^2$ to $\eta_k^R$. The squared tail and cutoff error is
$O(\rho^{2\nu+2})$.

Let $h_\epsilon^R(k)$ be the harmonic representative on the model of the
corresponding surviving class. Extend \eqref{eq:closed-compact-approximation}
by zero over the caps. Minimization first in the fixed model and then
$\rho\downarrow0$ gives
\[
 \limsup_{\epsilon\to0}\|h_\epsilon^R(k)\|^2\le G_R(k,k).
\]
For nonzero $k$, a fixed compactly supported closed dual four-form bounds
these norms below. Proposition~\ref{prop:multicap-compactness} at eigenvalue
zero gives a central limit without lost norm. Compactly supported dual
pairings identify its class as $k$, and Lemma~\ref{lem:central-kernels}
identifies the limit as $\eta_k^R$. Consequently the full harmonic Gram
matrix on $K$ converges to $G_R|_K$. All elliptic convergence used here
is on the fixed central model metric. No assertion about derivatives of
harmonic representatives for arbitrary metrics satisfying \textup{(N)}
is needed.

\subsection{Upper and lower bounds on every model subspace}
Let $D\in K^{\perp_{G_R}}$. On end $i$ write
\[
 \eta_D^R=r_i(D)\vartheta_i+dv_{+,D},\qquad
 |dv_{+,D}|=O(r^{\kappa-2}),\quad \kappa>0.
\]
Take the Lipschitz cutoff equal to zero below $2A\epsilon$, equal to one
above $R$, and equal to
\[
 \chi_\epsilon(r)=
 \frac{\log(r/(2A\epsilon))}{\log(R/(2A\epsilon))}
\]
between these radii. The form $T_\epsilon D=\chi_\epsilon\eta_D^R$,
extended by zero over the caps, lies in the model Hodge form domain.
The metric is exactly conical where the cutoff varies. Closedness and
coclosedness of $\eta_D^R$ give
$dT_\epsilon D=d\chi_\epsilon\wedge\eta_D^R$ and
$d^*T_\epsilon D=-\iota_{\nabla\chi_\epsilon}\eta_D^R$.
The radial contraction of $\vartheta_i$ vanishes, and its tangential
pairing with $d_Lv_{+,D}$ is zero. With $L_\epsilon=\log(1/\epsilon)$
and $c=2\pi/3$, it follows that
\begin{align}
 \bar q_\epsilon(T_\epsilon D)
 &=\frac{c\mathcal E([D],[D])}{L_\epsilon}
       +O_{R,A}(L_\epsilon^{-2})\|D\|^2,
       \label{eq:matrix-trial-energy}\\
 \langle T_\epsilon D,T_\epsilon F\rangle
 &=\mathcal M_R([D],[F])+O_{R,A}(L_\epsilon^{-1})\|D\|\|F\|.
       \label{eq:matrix-trial-mass}
\end{align}
Here any fixed norm on $V$ may be used. The remainder energy and mass
error are bounded, respectively, by
\[
 CL_\epsilon^{-2}\int_0^R r^{2\kappa-1}\,dr,
 \qquad CL_\epsilon^{-1}\int_0^R r(1+|\log r|)\,dr+O(\epsilon^2).
\]
Polarization gives the bilinear estimates, uniformly on this fixed
finite-dimensional space.

Subtract the full model harmonic projection. Its norm is $o(1)$ by the
preceding harmonic convergence and $G_R(D,K)=0$; the energy is unchanged.
Min--max on the first $j$ generalized eigendirections of $(\mathcal E,\mathcal M_R)$ gives
\begin{equation}\label{eq:matrix-upper}
 \limsup_{\epsilon\to0}L_\epsilon\bar\lambda_j(\epsilon;R,A)
 \le c\mu_j(R),\qquad 1\le j\le m.
\end{equation}
The $o(1)$ projection estimate suffices for this leading limit and does
not assert a second-order eigenvalue remainder.

For the lower bound choose a sequence on which
$L_\epsilon\bar\lambda_j$ tends to its lower limit $\ell$; the upper
bound makes it finite. The first $j$ orthonormal model eigenforms have a
common subsequence with orthonormal central limits
$v_1,\ldots,v_j\in K^{\perp_{G_R}}$, by
Proposition~\ref{prop:multicap-compactness} and convergence of the model
kernel. For $z\in\RR^j$ their combination $u_\epsilon(z)$ has
\[
 \bar q_\epsilon(u_\epsilon(z))\le\bar\lambda_j|z|^2,
 \qquad v(z)=\sum_{a=1}^jz_av_a.
\]
At the fixed sections $r_*=R/2$, its harmonic boundary coefficients
converge smoothly to $r_i(v(z))$. Each inner region is a scaled fixed
cap with an exact conical end. Lemma~\ref{lem:capacity} gives
\[
 \bar q_\epsilon(u_\epsilon(z))\ge
 \sum_i\frac{c|B_{i,\epsilon}(z)|^2}
                    {\log(r_*/\epsilon)+C_{i,A}}.
\]
The regions are disjoint. Multiplying by $L_\epsilon$ and taking the
limit yields, simultaneously for the finite basis and hence every $z$,
\[
 c\mathcal E(v(z),v(z))\le\ell|z|^2=\ell \mathcal M_R(v(z),v(z)).
\]
The largest generalized Rayleigh quotient on this $j$-dimensional span is
at least $\mu_j(R)$. Together with \eqref{eq:matrix-upper}, this proves
\begin{equation}\label{eq:matrix-model-law}
 \lim_{\epsilon\to0}L_\epsilon\bar\lambda_j(\epsilon;R,A)
       =c\mu_j(R),\qquad 1\le j\le m.
\end{equation}
The constants $C_{i,A}$ require no uniform bound as $A$ varies, since the
family limit has been taken first.

\subsection{Transfer under \texorpdfstring{$C^0$}{C0} comparison}
\begin{lemma}[Dodziuk comparison]\label{lem:metric-spectra}
Let $g,h$ be smooth metrics on a closed $n$-manifold and suppose
$(1-\delta)g\le h\le(1+\delta)g$, with $0<\delta<1/2$.
For every ordered positive Hodge eigenvalue in degree $p$,
\[
 K(\delta)^{-1}\lambda_{p,j}^+(g)
 \le\lambda_{p,j}^+(h)\le K(\delta)\lambda_{p,j}^+(g),
 \qquad K(\delta)=1+O_n(\delta).
\]
No derivatives of the metric comparison are required.
\end{lemma}
\begin{proof}
This is the closed-manifold form of \cite[Proposition~3.3]{dodziuk}.
We recall the variational argument to specify the metric dependence.
The squared $L^2$ norms of $p$-forms satisfy
\[
 a_p\|u\|_g^2\le\|u\|_h^2\le b_p\|u\|_g^2,
 \quad a_p=(1-\delta)^{n/2}(1+\delta)^{-p},\quad
 b_p=(1+\delta)^{n/2}(1-\delta)^{-p}.
\]
The distributional differential $d$ and its graph domain are unchanged
under equivalent norms. Set $Z^p=\ker d$ and give
$\operatorname{Dom}d/Z^p$ the quotient norm
$\|[u]\|_{g,q}^2=\inf_{z\in Z^p}\|u+z\|_g^2$.
The same constants compare these quotient norms. The coexact positive
$p$-eigenvalues are the squared singular values of $d$ on this quotient,
with min--max quotient $\|du\|_{p+1,g}^2/\|[u]\|_{g,q}^2$.
Its comparison factors are $a_{p+1}/b_p$ and $b_{p+1}/a_p$.
The positive exact $p$-spectrum equals the coexact $(p-1)$-spectrum.
The full positive spectrum is the ordered union of these two lists,
with multiplicity; the same common multiplicative bounds hold for their
ordered union. This proves the lemma and avoids comparing codifferentials
pointwise.
\end{proof}

\begin{theorem}[Spectral consequence of condition \textup{(N)}]
\label{thm:C0-spectrum}
Under the geometric hypotheses of Theorem~\ref{thm:matrix}, let
$h_\epsilon$ be any smooth metric family satisfying condition \textup{(N)}
relative to the central metric $g_0$. Then exactly $m$ positive degree-two
eigenvalues tend to zero, with
\[
 \lim_{\epsilon\downarrow0}\log(1/\epsilon)
             \lambda^+_{2,j}(h_\epsilon)=\frac{2\pi}{3}\mu_j,
 \qquad 1\le j\le m.
\]
There are uniform positive complementary gaps in degrees one and two.
If these metrics are Ricci-flat K\"ahler with trivial canonical bundle,
the small spectral space is coexact and primitive of real type $(1,1)$.
\end{theorem}
\begin{proof}
Fix $R,A$. For every $\eta>0$ and all sufficiently small $\epsilon$,
condition \textup{(N)} bounds the relative metric error by
$\delta(R,A)+\eta$. Take this sum small enough for
Lemma~\ref{lem:metric-spectra}. Its multiplicative factor
$K(\delta+\eta)$ bounds each ordered positive eigenvalue above and below.
Using \eqref{eq:matrix-model-law}, then letting $\eta\downarrow0$,
$R\downarrow0$ and $A\to\infty$, proves the coefficient because
$\mu_j(R)\to\mu_j$. One fixed admissible pair $R,A$ transfers the
model complementary gaps by a fixed positive factor. No gap uniform in
the auxiliary radii is needed. The K\"ahler assertion follows from the
next lemma.
\end{proof}

\begin{lemma}\label{lem:small-type}
On a compact Ricci-flat K\"ahler threefold with trivial canonical bundle,
every real two-form eigenvector with positive eigenvalue below the first
positive one-form eigenvalue is coexact and primitive of type $(1,1)$.
\end{lemma}
\begin{proof}
A nowhere-zero holomorphic three-form $\Omega$ is parallel. Using the
complex-bilinear metric to raise indices, the map
\[
 T:\Lambda^{0,1}\longrightarrow\Lambda^{2,0},\qquad
 T\beta=\iota_{\beta^\sharp}\Omega
\]
is a complex-linear parallel bundle isomorphism, equivariant for $SU(3)$
and hence for the actual holonomy group. It intertwines the Hodge
Laplacians: both the rough term and the curvature term $q(R)$ in the
Weitzenb\"ock formula are natural under this holonomy intertwiner
\cite[\S\S2.4, 2.6, Definition~2.4 and Example~1]{semmelmann-weingart}.
Explicitly, $q(R)=\sum_\alpha\rho(X_\alpha)\rho(R(X_\alpha))$ in an
orthonormal holonomy-algebra frame. Thus a nonzero $(2,0)$ eigencomponent
would give a nonzero $(0,1)$ eigenform of the same eigenvalue via $T^{-1}$,
contradicting the one-form gap. Reality removes the conjugate component.

K\"ahler commutation makes $\Lambda u$ a scalar eigenfunction of the
same positive eigenvalue. If nonzero, its differential gives the same
contradiction. Similarly $d^*u$ must vanish. Hence $\Lambda u=0$ and
$u=\lambda^{-1}d^*du$. The argument applies to the entire eigenspace and
requires no simplicity hypothesis.
\end{proof}

Lemma~\ref{lem:model-comparison} establishes condition \textup{(N)} for
the actual polarized Ricci-flat metrics. Substituting
$\epsilon=t^{1/3}$ in Theorem~\ref{thm:C0-spectrum} proves
Theorem~\ref{thm:matrix}. A simple $\mu_j$ is separated from its neighbours;
the ordered laws and complete count imply eventual simplicity. In rank
one the quotient mass is \eqref{eq:mass-identity}, so the same proof gives
all assertions of Theorem~\ref{thm:main} except its sharper upper remainder.

\subsection{The differentiated rank-one remainder}
The remaining assertion uses the stronger full-central estimate
\eqref{eq:actual-comparison}, including its first derivative. Impose
$b_2(Y)=2$, $b_2(X_t)=1$ for the rest of this subsection.
We now use the global form $\alpha$ of Proposition~\ref{prop:mass}.
Put $L_\epsilon=\log(1/\epsilon)$. On end $i$ choose a cutoff
$\chi_{i,\epsilon}$ vanishing below $A_i\epsilon$, equal to one
above a fixed small $R_i$, and logarithmic between these radii:
\[
 \chi_{i,\epsilon}(r)=
 \frac{\log(r/(A_i\epsilon))}{\log(R_i/(A_i\epsilon))}.
\]
Round its endpoints smoothly on fixed multiplicative annuli, making
it flat at the inner endpoint. The rounding changes energy by
$O(L_\epsilon^{-2})$. Denote the denominator by
$L_i=L_\epsilon+O(1)$. Transport $\chi_\epsilon\alpha$ using
$F_\epsilon$ and extend by zero across the cores, obtaining a smooth
real two-form $u_\epsilon$ on $X_t$.

\begin{proposition}\label{prop:upper}
For the actual Ricci-flat metrics,
\begin{equation}\label{eq:epsilon-upper}
 \lambda_2^+(\epsilon)
 \le\frac{c\sum_i b_i^2}{M L_\epsilon}+O(L_\epsilon^{-2}),
 \qquad c=2\pi/3.
\end{equation}
\end{proposition}

\begin{proof}
Pullback commutes with $d$, so $du_\epsilon$ is the transport of
$d\chi_\epsilon\wedge\alpha$. On end $i$ its leading energy is
$c b_i^2\int r|\chi_i'|^2dr=c b_i^2/L_i+O(L_\epsilon^{-2})$.
The errors in the form, central metric and actual metric carry respectively
positive factors $r^\nu$, $r^\mu$ and $(\epsilon/r)^{a_0}$.
Their integrals against $dr/r$ are bounded, and each has prefactor
$O(L_\epsilon^{-2})$. Thus
\begin{equation}\label{eq:d-trial}
 \|du_\epsilon\|^2=\frac{c\sum_i b_i^2}{L_\epsilon}
 +O(L_\epsilon^{-2}).
\end{equation}

The actual codifferential requires a separate estimate.
Proposition~\ref{prop:metric-comparison} and the differentiated
$O(r^{-2})$ bound for $\alpha$ give
\[
 |(d^*_{g_\epsilon}-d^*_{g_0})\alpha|
 \le C(\epsilon/r)^{a_0}r^{-3}.
\]
Since $d^*_{g_0}\alpha=0$, the contribution multiplied by $\chi_i$
has squared norm at most
\[
 C\int_{A_i\epsilon}^{R_i}\chi_i^2(\epsilon/r)^{2a_0}\frac{dr}{r}
 \le \frac C{L_\epsilon^2}
 \int_{A_i}^{\infty}(1+\log(s/A_i))^2s^{-1-2a_0}ds
 =O(L_\epsilon^{-2}).
\]
The product-rule contraction obeys
\[
 |\iota_{\nabla_{g_\epsilon}\chi_i}\alpha|
 \le C|\chi_i'|\bigl(r^{\nu-2}+(\epsilon/r)^{a_0}r^{-2}\bigr),
\]
because the radial cone contraction of $\vartheta_i$ vanishes.
Its squared norm has bounded radial integrals
$\int r^{2\nu-1}dr$ and $\int(\epsilon/r)^{2a_0}dr/r$, again with
prefactor $O(L_\epsilon^{-2})$. On the fixed exterior the squared
error is $O(\epsilon^{2a_0})$. Combining with \eqref{eq:d-trial},
\begin{equation}\label{eq:q-trial}
 q_\epsilon(u_\epsilon):=\|du_\epsilon\|^2+\|d^*u_\epsilon\|^2
 =\frac{c\sum_i b_i^2}{L_\epsilon}+O(L_\epsilon^{-2}).
\end{equation}

The mass density of $\alpha$ is $O(r\,dr)$, and the removed inner
mass is $O(\epsilon^2)$. On the remainder,
$|1-\chi_i|\le C(1+|\log r|)/L_\epsilon$, which is integrable
against that density. The actual metric error integrates to
$O(\epsilon^{\min(a_0,2)}(1+|\log\epsilon|))$. Therefore
\[
 \|u_\epsilon\|^2=M+O(L_\epsilon^{-1}).
\]
The K\"ahler-form comparison and pointwise central primitivity similarly
give $\langle u_\epsilon,\omega_\epsilon\rangle=O(L_\epsilon^{-1})$.
Subtract its orthogonal projection onto $\omega_\epsilon$.
Since $\|\omega_\epsilon\|^2=v/2$, this changes the squared norm by
$O(L_\epsilon^{-2})$ and leaves $q_\epsilon$ unchanged. It removes
the whole harmonic kernel because $b_2(X_t)=1$. The Rayleigh principle
then proves \eqref{eq:epsilon-upper}.
\end{proof}

Since $\log(1/\epsilon)=|\log t|/3$, Proposition~\ref{prop:upper}
gives \eqref{eq:main-upper} and completes the proof of
Theorem~\ref{thm:main}. This estimate is separate from the $C^0$ transfer
and has only the stated one-sided conclusion.
